% Portable, self-contained Pandoc publication template.
\documentclass[11pt,letterpaper]{article}
\usepackage[T1]{fontenc}
\usepackage{lmodern}
\usepackage{amsmath,amssymb,amsthm,mathtools}
\usepackage[margin=0.9in]{geometry}
\usepackage{microtype}
\usepackage{booktabs,longtable,array,calc,multirow}
\usepackage{etoolbox}
\usepackage{graphicx}
\usepackage{xurl}
\usepackage[numbers,sort&compress]{natbib}
\usepackage[hidelinks,unicode]{hyperref}
\usepackage{bookmark}
\urlstyle{same}
\setcounter{secnumdepth}{3}
\setlength{\emergencystretch}{3em}
\setlength{\parskip}{0.2em}
\allowdisplaybreaks[1]
\providecommand{\tightlist}{%
  \setlength{\itemsep}{0pt}\setlength{\parskip}{0pt}}
\providecommand{\passthrough}[1]{#1}
\providecommand{\pandocbounded}[1]{#1}
\AtBeginEnvironment{longtable}{\small\setlength{\tabcolsep}{4pt}}
\newenvironment{manualstatement}{%
  \par\addvspace{0.7\baselineskip}\begingroup\itshape\noindent\ignorespaces
}{\par\endgroup\addvspace{0.4\baselineskip}}
\newenvironment{manualdefinition}{%
  \par\addvspace{0.7\baselineskip}\begingroup\normalfont\noindent\ignorespaces
}{\par\endgroup\addvspace{0.4\baselineskip}}
\newenvironment{manualproof}{%
  \par\addvspace{0.4\baselineskip}\begingroup\normalfont\noindent\ignorespaces
}{\par\endgroup\addvspace{0.5\baselineskip}}
\title{Refining Hardness after ETH\\[0.5em]{\large Task-Relevant
Distinctions and the Cost of Exact Counting}}
\author{Rohan Vinaik}
\date{\small 7 November 2026}
\begin{document}
\maketitle
\begin{abstract}
The reported subexponential algorithms for satisfiability leave a strong
successor question: can exact counting gain a fixed power in its
exponent? We examine that question through the distinctions a task
requires, the operations that resolve them jointly, and what those
operations cost, with the public hypotheses \#SH and Weak \#ETH as the
quantitative setting. After comparing the existing method-specific
barriers, we give a fully formalized finite analysis of the
coordinate-swap construction behind one interpolation algorithm. We
construct the measurement rows it selects, prove their fixing-kernel
property and their exact row count, and connect them to a uniform
tradeoff among measurement, degree and enumeration. Under the stated row
and dimension budgets, the complete clipped substitution enumeration has
size at least \(2^{n\log_2\log_2 n/(4\sqrt{\log_2 n})}\) for
\(n\ge2^{16}\), and eventually at least \(2^{n^{1-\varepsilon}}\) for
every fixed \(\varepsilon>0\). This bounds the enumeration as completed;
an algorithm that evaluates its aggregate by other operations calls for
a separate analysis. The supporting mathematical implications and the
constructed-measurement application are checked in Lean 4, and the
unrestricted counting hypothesis stays open.
\end{abstract}
\section{Introduction}\label{introduction}

The Exponential Time Hypothesis fixes the difficulty of 3-SAT at one
scale: no deterministic algorithm runs in time \(2^{o(n)}\), up to
polynomial factors in the input, where \(n\) counts the variables
\citep{IPZ01}. Two recent preprints report algorithms that cross this
boundary \citep{O26, Ant26}. Both algorithms count satisfying
assignments exactly rather than merely deciding satisfiability. If their
analyses hold, ETH is false. What they leave open is how much further
the exponent can fall.

That gap is both quantitative and structural. A running time of
\(2^{n/\sqrt{\log n}}\) is subexponential, yet for every fixed
\(\varepsilon>0\) its exponent still exceeds \(n^{1-\varepsilon}\) once
\(n\) is large. A conjecture that rules out the latter scale therefore
survives such an algorithm untouched. The Satisfiability Hypothesis of
Stearns and Hunt \citep{SH90}, with its counting analogue, asks exactly
this successor question. MIT Hardness Group \citep{SH26} have since
shown it equivalent to Weak ETH and Weak \#ETH, and they have built the
matching reduction theory.

We ask what the distinction means for an account of hardness. A problem
is hard neither because its presentation has many variables nor because
its solution space has many candidates. One operation can resolve many
apparent alternatives at once, while a compact representation of the
answer can still be expensive to evaluate. So a refined hardness claim
has to name four things together: the task, the alternatives that change
its answer, the operations permitted on the input, and the cost of
obtaining and using what those operations return.

Our thesis is that the failure of ETH's original scale calls for this
more discriminating account, and not for abandoning strong hardness
hypotheses. Exact counting is a sharp setting for it: its answer is
short, but computing that answer means controlling interactions among
many assignments. Both reported algorithms show that those interactions
can be organized more cheaply than ETH allowed. What remains is whether
successive reorganizations can lower the exponent by a fixed power, or
whether a near-exponential obstruction persists. We take \#SH, already
public, as the quantitative form of that question, and we introduce no
new name for it.

Sections 3 and 4.1 mark two limits on what an information argument alone
can prove. Section 3 gives an exponential point-query separation between
counting and decision on small CNFs. Reading the formula as given costs
only its encoding length in bit queries, so carrying that separation to
running time would need a justified access simulation and its own
account of computation. Section 4.1 gives the exact Hermite threshold
for arbitrary degree-bounded polynomials. Its necessity proof uses
indistinguishable polynomials that need not arise from formulas, and its
threshold counts coefficients supplied rather than computed, so it says
nothing yet about the cost of producing them.

Those two checks bound the lower-bound argument from the outside;
Theorem 5.1 rests on neither. Section 4 looks at the work the algorithms
do on the input as written. Section 5 then finds an obstruction inside
the operations of one such construction. Its coordinate swaps have an
exact spectrum, so the fixing condition imposes a rank requirement on
every rational measurement map. Agreement with parity turns that
requirement into a degree floor, and the degree and dimension budgets
then force a near-exponential complete substitution enumeration. We also
build the measurement rows that the source algorithm selects, and we
prove that they satisfy the fixing property.

The spectral mechanism and the near-exponential template scale both have
public antecedents, which Section 4.3 discusses. We add the finite,
uniformly quantified implication itself, its mechanization on the
coordinate space the algorithm uses, and the checked application to the
source's own measurements. Together these give one concrete instance of
the broader structural program: a result that says which changes of
fixing polynomial or measurement rows yield no fixed-power improvement
while the other contracts and the complete enumeration stay in place.
Whether the unrestricted successor hypothesis holds is another question,
and it stays open.

\section{The quantitative question after
ETH}\label{the-quantitative-question-after-eth}

\subsection{Exponential scale and power of the
exponent}\label{exponential-scale-and-power-of-the-exponent}

Let \(L\) denote a structural encoding length of a CNF formula, counting
its variables, clauses and literal occurrences. Conventional indexed
binary encodings change this length by at most polylogarithmic factors,
and such changes preserve the fixed-power separation in the hypotheses
below. An arbitrary polynomial blowup needs its own accounting.
Following the conventions of \citep{SH26}, we include unused variables
in \(L\), and the variable-parameter hypothesis keeps polynomial
running-time factors in the full input length in view.

\textbf{Hypothesis 2.1 (Weak \#ETH and \#SH).} Weak \#ETH asserts that
for every fixed \(\varepsilon>0\), no deterministic algorithm computes
the exact number of satisfying assignments of a 3-CNF formula in time
\[O\!\left(2^{n^{1-\varepsilon}}\right)L^{O(1)}.\] The hypothesis \#SH
asserts that, for every fixed \(\varepsilon>0\), exact counting for
unrestricted CNF formulas has no deterministic running time
\(O(2^{L^{1-\varepsilon}})\).

MIT Hardness Group prove that these two hypotheses are equivalent
\citep[Theorem 2.4]{SH26}. Their counting-preserving ingredients are a
linear-size Tseitin conversion with unique gate extensions and a
disjoint, quantitative sparsification procedure, and their decision
versions relate SH and Weak ETH in the same way. Both are public
hypotheses with public reductions, so our unconditional theorems assume
neither.

This distinction from ETH is more than terminology. For any fixed
\(a>0\) and fixed real \(b\), a positive exponent of order
\[\frac{n(\log\log n)^b}{(\log n)^a}\] is \(o(n)\) and yet
\(n^{1-o(1)}\). An algorithm at such a scale therefore refutes ETH while
remaining compatible with Weak \#ETH. Calling it subexponential is
correct; concluding that it lowers the exponent by a fixed power would
be wrong.

Both source results sit precisely there. For each fixed clause width,
OpenAI report a counting time on general formulas of that width of
\[L^{O(1)}\exp\!\left(O\!\left(\frac{n\log\log n}{\sqrt{\log n}}\right)\right)\]
and a sparse-formula bound with exponent \(O(n/\sqrt{\log n})\)
\citep{O26}. The Claude paper reports, for fixed clause width \(k\), the
bound \[2^{\,O\!\left(\min\left\{
\frac{km}{\sqrt{\log m}},
\frac{k^2 n\log\log n}{\sqrt{\log n}}
\right\}\right)},\] where \(m\) counts clauses and \(n\) counts the
variables that occur \citep{Ant26}. Unused variables contribute a factor
to the count that is computed directly. We take both as the upper bounds
their authors state; the present formalization has not re-verified them.

\subsection{What a structural explanation must
add}\label{what-a-structural-explanation-must-add}

\#SH states a rate of growth, and nothing more. It leaves unexplained
which aspect of the input prevents a faster computation. An account of
that obstruction needs more than the number of possible assignments,
because summation, elimination, symmetry and algebraic identities all
treat many assignments together.

We therefore ask three questions. Which alternative instances require
different task answers? Which permitted operations distinguish those
alternatives? And how much computation goes into selecting,
constructing, evaluating and combining those operations? Two of these
fix an information requirement, while the third asks how an algorithm
realizes it, and the analysis below needs all three.

This separation also clarifies what a change of representation can do. A
new representation may preserve every task-relevant distinction while
making their joint resolution far cheaper. It may instead introduce an
operation that is substantially more informative than the old ones. A
certificate bound transfers only when a stated simulation preserves the
relevant access costs, and a running-time claim needs computational
control of the representation and its decoder as well. So there is no
general inference from preserving information to preserving
computational difficulty.

Our structural reading of \#SH is correspondingly specific. Can every
attempt at a fixed-power improvement be shown to leave an expensive
unresolved interaction, once all jointly computable operations have been
charged? Posed for unrestricted algorithms that is a research question.
Section 5 answers one instance of it inside a defined construction, with
the operators the algorithm uses and with stated budgets. Section 6 sets
out what still stands between that instance and an unrestricted result.

\section{Distinctions required by a
task}\label{distinctions-required-by-a-task}

\subsection{Certificates and completed
computations}\label{certificates-and-completed-computations}

Let \(\mathcal C\) be a class of admissible instances, let
\(q:\mathcal C\to Y\) be the task, and let each \(e\in E\) specify an
observation \(o_e:\mathcal C\to A\). The answer type may encode several
scalars or a structured object. The class \(\mathcal C\) belongs to the
model, since it says which alternatives a computation has to handle.

\begin{manualdefinition}

{{\upshape \textbf{Definition 3.1 (Task certificate).}}} A finite set
\(T\subseteq E\) is a certificate for an instance \(x\in\mathcal C\) if
\[\forall y\in\mathcal C,\quad
\bigl(\forall e\in T,\ o_e(y)=o_e(x)\bigr)
\Longrightarrow q(y)=q(x).\] Write \(\operatorname{cert}_q(x)\) for its
minimum cardinality, with value \(\infty\) when no finite certificate
exists. With nonnegative observation charges \(c_e\), use the infimum of
\(\sum_{e\in T}c_e\) over finite certificates instead, again taking
\(\infty\) for an empty family.

\end{manualdefinition}

This is the usual certificate notion for the declared observations,
extended to a general task, and we claim no new complexity measure for
algorithms in it: standard decision-tree and priced-information
treatments already distinguish certificates, queries and the cost of
finding useful information \citetext{\citealp[Section
4.1]{BW02}; \citealp{CFGKRS02}}.

A deterministic observation policy chooses its next observation, or its
final answer, from the previous query--answer history alone. Any further
access to the instance therefore has to appear among the permitted
observations.

\begin{manualstatement}

{{\upshape \textbf{Proposition 3.2 (Transcript certificate).}}} If a
deterministic policy halts on \(x\) with the correct task answer and is
correct whenever it halts on an admissible instance, its completed
observation transcript is a certificate for \(x\). In particular, the
number of observation calls on a correct terminating run is at least
\(\operatorname{cert}_q(x)\).

\end{manualstatement}

\begin{manualproof}

\textbf{Proof.} An instance agreeing with the completed transcript
produces the same first answer, hence the same next query. By induction
it reproduces every query and answer and the same stopping decision, and
correctness then gives the same task value. Removing repeated
observations can only shrink the transcript. \(\square\)

\end{manualproof}

An elementary capacity bound shows the role of task-relevant
alternatives directly. Suppose \(W\subseteq\mathcal C\) consists of
\(M\) instances with \(q(y)\ne q(x)\), and every observation disagrees
with \(x\) on at most \(r\) members of \(W\). Every finite certificate
\(T\) must exclude all of \(W\); hence \[M\le r\,|T|.\] The proof is a
union bound over the alternatives each observation excludes. For \(r>0\)
it gives \(\operatorname{cert}_q(x)\ge M/r\), infinite certificate
complexity included. A large \(M\) on its own proves nothing when one
observation can exclude all \(M\) alternatives; the capacity of the
permitted operation is essential.

\subsection{Counting and decision can require different
distinctions}\label{counting-and-decision-can-require-different-distinctions}

Let \(X=\{0,1\}^n\), \(n\ge1\), and take as admissible instances the
Boolean functions represented by CNFs of width at most two. Observations
reveal the function value at one assignment, and the tasks are
satisfiability and exact count. The alternative formulas used in the
proof below have at most \(2n\) clauses.

\begin{manualstatement}

{{\upshape \textbf{Proposition 3.3 (Point-access separation on small
CNFs).}}} For every \(a\in X\), the function with satisfying set
\(\{a\}\) has a one-observation satisfiability certificate and an
exact-count certificate of size \(2^n\) in this model. All \(2^n\) point
observations suffice for counting any admissible instance.

\end{manualstatement}

\begin{manualproof}

\textbf{Proof.} A positive answer at \(a\) certifies satisfiability. The
empty transcript certifies nothing, since the identically false function
is admissible.

For counting, the point \(a\) must be observed to separate \(\{a\}\)
from the empty satisfying set, and every other point \(b\) must be
observed to separate \(\{a\}\) from \(\{a,b\}\). Both alternatives lie
in the stated CNF class: a singleton is specified by unit clauses, and
the false function by two contradictory units. To specify \(\{a,b\}\)
when \(a\ne b\), fix all coordinates on which they agree, choose a
coordinate on which they differ, and require each other differing
coordinate to follow it with the appropriate sign. Each such equivalence
uses two width-two clauses, so the total is at most \(2n\). Thus all
points are necessary, and observing all of them is sufficient.
\(\square\)

\end{manualproof}

The example isolates a difference between tasks; it is no hard family
for 2-CNF formulas supplied in full. Under point access, counting cannot
stop at a satisfying witness, because an unobserved assignment may
change the count. A written clause list gives different information and
supports different operations. Indeed, the two-point formulas in the
proof expose their small satisfying sets on sight.

For an \(L\)-bit input given in full, reading the \(L\) bits
reconstructs the entire instance. Every task then has an acquisition
procedure using at most \(L\) bit observations, once subsequent internal
computation goes uncharged. We formalize this operational statement as
well. It rules out using a point-access information bound as a
running-time lower bound for arbitrary computation on a formula given in
full.

\subsection{A sufficient condition for transfer between
representations}\label{a-sufficient-condition-for-transfer-between-representations}

Suppose a representation map \(R:\mathcal C\to\mathcal C'\) preserves
admissibility and the task: \(q'(R(x))=q(x)\). For each new observation
\(e'\), suppose a specified finite set \(S(e')\subseteq E\) and a
decoder satisfy \[o'_{e'}(R(x))=
 d_{e'}\bigl((o_e(x))_{e\in S(e')}\bigr)
\qquad(x\in\mathcal C).\] The sets are fixed for each \(e'\), so this is
a finite, nonadaptive simulation of an individual observation.

\begin{manualstatement}

{{\upshape \textbf{Proposition 3.4 (Observation simulation).}}} Assume
the new observation family is finite. If \(K>0\) and
\[\sum_{e\in S(e')}c_e\le Kc'_{e'}
\quad\text{for every }e',\] then the weighted certificate infima satisfy
\[\operatorname{cert}_{q,c}(x)
\le K\operatorname{cert}_{q',c'}(R(x)).\]

\end{manualstatement}

\begin{manualproof}

\textbf{Proof.} Replace each observation in a new certificate \(T'\) by
\(S(e')\). Agreement on their union gives equal decoded answers, so task
preservation transports the certificate to \(x\). Nonnegative costs
bound the union's cost by the sum of its parts, at most
\(K\sum_{e'\in T'}c'_{e'}\). Since the new observation family is finite,
an attaining certificate exists whenever any certificate does, which
gives the infimum comparison. A missing new certificate gives the value
\(\infty\). \(\square\)

\end{manualproof}

The proposition compares acquisition costs. Its decoder is a
mathematical function, so it carries no efficiency guarantee.
Representation construction, decoding, arithmetic, query selection and
the reuse of intermediate results each need their own running-time
analysis. In particular, a missing simulation leaves a faster algorithm
unrefuted; it only means that this transfer argument is unavailable.

\section{What the algebraic algorithms
compute}\label{what-the-algebraic-algorithms-compute}

\subsection{Recovery from polynomial
data}\label{recovery-from-polynomial-data}

In the rational interpolation route of \citep[Sections 2.4--2.7]{O26}, a
bounded-occurrence formula is encoded by a univariate polynomial \(H\)
with \[H(0)=\#F,\qquad \deg H\le MR\le3nR.\] Here \(M\) is a
total-incidence budget and \(R\) controls the fixing polynomial. The
algorithm computes initial Taylor coefficients at \(Q\) interpolation
points and recovers the desired value from them. For arbitrary
bounded-degree polynomials, the information threshold behind that
recovery is exact.

\begin{manualstatement}

{{\upshape \textbf{Proposition 4.1 (Ambient Hermite threshold).}}} Let
\(t_1,\ldots,t_Q\) be distinct nonzero rational numbers, let
\(J_i\in\mathbb N\), and let \(d\in\mathbb N\). Among rational
polynomials of degree at most \(d\), the coefficients of \((X-t_i)^k\)
for integers \(0\le k<J_i\) determine the value at zero for every
polynomial if and only if \[\sum_{i=1}^Q J_i>d.\]

\end{manualstatement}

\begin{manualproof}

\textbf{Proof.} Equality of the jets makes the polynomial difference
divisible by \[P(X)=\prod_{i=1}^Q(X-t_i)^{J_i}.\] Distinct nodes give
coprime factors. If \(\deg P>d\), divisibility forces the difference to
vanish. If \(\deg P\le d\), then \(H\) and \(H+P\) have the same
requested jets and stay within the degree bound, yet they differ at zero
because \(P(0)\ne0\). \(\square\)

\end{manualproof}

Here the ambient class matters, because adding \(P\) need not produce
another polynomial generated by a CNF encoding. Its necessity direction
therefore leaves open whether the threshold is optimal for that
restricted family, or for every way of recovering its count. For
\(Q\ge1\), sufficiency does justify the standard choice
\[q=\left\lfloor\frac{3nR}{Q}\right\rfloor,
\qquad J=q+1,
\qquad QJ>3nR.\]

The computational question is how the jets are obtained. In the cited
construction, shared integer measurements separate variable groups.
Taylor coefficients of order at most \(q\) then reduce to small-support
substitutions and group sums, and the complete substitution family has
size \[A(n,q)=\sum_{j=0}^{\min(q,n)}2^j\binom nj.\] In the source's
accounting this factor multiplies the enumeration of measurement values
and of assignments inside a group, giving the representative product
\[A(n,q)\,(2D+1)^{2\tau n}\,2^{2n/r},\] up to further polynomial and
implementation factors \citep{O26}. To evaluate the finite-difference
identity, the routine visits every pair \(U\subseteq E\subseteq[n]\)
with \(|E|\le q\). Here \(U\) fixes the Taylor correction, and \(E\)
fixes which assignment coordinates must equal one in the inner sum. For
each pair the routine scans shared measurement tuples and forms products
of local group sums; one correction with one measurement tuple can touch
several groups \citep[Sections 3.3--3.5]{O26}. Faster interpolation of
the resulting jets would therefore leave this loop in place. Avoiding it
means computing the signed aggregate by some other procedure, and the
dependence across groups still allows reuse, dynamic programming or a
different factorization.

\subsection{Exact network identities and their
implementation}\label{exact-network-identities-and-their-implementation}

The Claude algorithm builds a network with a bit for each literal
occurrence, clause factors, and equality factors linking the occurrences
of one variable. Its network sum is the exact count. An exact crossing
identity replaces a crossing by a signed sum of pieces with a simpler
row or column direction, and grouping and recursion then compute the
resulting sums \citep[Sections 3--6]{Ant26}.

Its abstract network model charges for factor and membership evaluation
and bounds integer magnitudes. In the implementation, a grouped table
query is computed by summing over the grouped local assignments rather
than read from a materialized table of free answers. Recursive calls
receive evaluator procedures, and the stronger variable bound also uses
truncated inclusion--exclusion and deterministic restrictions
\citep[Appendices B--D]{Ant26}.

This is a second way to resolve interactions jointly. It leaves the
number of assignments the answer concerns where it was; what it reduces
is the work of combining their contributions. Proposition 3.4 therefore
gives no identification of this network computation with point
observations, and the swap hypotheses of Section 5 have not been derived
for it.

Brito's subsequent analysis gives a related example of a
construction-specific restriction \citep{Brito26}. For the relevant
graded shape, rank subadditivity bounds the number of crossing
directions that its pieces can supply. In the application using all
basis labels, the crossing has rank \((\dim V_h)^2\), while the
directional pieces have total rank at most \((4h+1)\dim V_{\le4h}\).
This growth constraint concerns that shape and its dimension accounting
rather than counting in general. Our result below uses different
operators, and it proves their rank identity exactly.

\subsection{Related method-specific
obstructions}\label{related-method-specific-obstructions}

Dias and Lima study clause-product interpolation with restored
occurrence polynomials, arbitrary nonzero rational evaluation points,
and possibly unequal Taylor orders \citep[Sections 5--9]{DL26}. In their
setting, recovery must be exact for every polynomial up to the declared
degree bound. Their evaluation model charges for group size, interaction
width and Taylor order, with a further measurement-rank charge for
product-type evaluation. A vanishing-product argument ties these
parameters to the input's graph structure, which yields a
treewidth-scale obstruction and, on specified random formulas, a
\(2^{\Omega(\sqrt n)}\) model-cost bound. Their stated charges exclude
some alternative computations inside groups, and their scope discussion
explains why the sharper parity and tensor routes escape the model.
Theorem 5.1 instead fixes one swap operator family together with its
dimension budget. Its larger exponent concerns a more specific
enumeration contract, so the exponents alone decide nothing about which
result is the stronger one.

Majdoddin studies an arithmetic-cost class of rewrite-and-merge
algorithms for weighted clique sums \citep[Definition 4.2 and Section
4]{Maj26}. His model evaluates every product choice and charges
materialized tables and recursive calls. Its balanced homogeneous
two-merge subclass has a rate floor of \((\log_2K)^{-3}\) under the
stated base-method condition. With network size \(S=\Omega(n)\) and
\(K\) bounded by a polynomially bounded clause count, this excludes
fixed-power improvements. Stronger grid bounds and separate single-block
statements carry their own hypotheses. These are public
near-exponential-scale results for specified network recursions. They
supply none of the fixing, row or dimension conditions of our swap
theorem, and our arbitrary-measurement quantifier in turn leaves that
recursive algorithm class uncovered.

Tulsyan gives the closest comparison \citetext{\citealp[Section
5]{T26}; \citealp[Sections 1--3]{T26note}}. His supplementary note
records the defect-rank identity for the commutative swap algebra, and a
triangular-submatrix argument then bounds defect rank for a broader
support-constrained matrix class, noncommuting operators included. His
paper states an asymptotic near-exponential template-cost consequence on
linked bounded-occurrence formulas, using the Arora--Rao--Vazirani
structure theorem, the specified Taylor enumeration and a vanishing
measurement rate. Its full-label tail argument and its linked-input
argument rest on different assumptions. These claims already place the
near-exponential scale in this line of work. Our finite bound below
reaches it by a direct parity-agreement argument, an explicit dimension
budget and the complete clipped sum, so it needs neither a linked-input
hypothesis nor a geometric structure theorem. We formalize the entire
implication together with the measurement map, rather than take the
external template proposition as a premise.

The comparisons are summarized by the quantity each result constrains:

{\def\LTcaptype{none} % do not increment counter
\begin{longtable}[]{@{}
  >{\raggedright\arraybackslash}p{(\linewidth - 4\tabcolsep) * \real{0.2590}}
  >{\raggedright\arraybackslash}p{(\linewidth - 4\tabcolsep) * \real{0.4037}}
  >{\raggedright\arraybackslash}p{(\linewidth - 4\tabcolsep) * \real{0.3374}}@{}}
\toprule\noalign{}
\begin{minipage}[b]{\linewidth}\raggedright
Result
\end{minipage} & \begin{minipage}[b]{\linewidth}\raggedright
Permitted mechanism
\end{minipage} & \begin{minipage}[b]{\linewidth}\raggedright
Bounded quantity
\end{minipage} \\
\midrule\noalign{}
\endhead
\bottomrule\noalign{}
\endlastfoot
Dias--Lima & Clause-product schemes with declared evaluation charges &
Model evaluation cost \\
Majdoddin & Restricted recursive rewrites and merges & Arithmetic-cost
exponent \\
Tulsyan & Support-constrained template matrices and specified input
conditions & Defect rank and asymptotic template cost \\
Theorem 5.1 & Full swap history, univariate fixing, row and dimension
budgets & Complete clipped enumeration cardinality \\
\end{longtable}
}

An implication between these models needs a construction that preserves
their hypotheses and charges. In particular, the freedom to choose any
rational measurement map for a fixed operator yields a bound for that
operator alone. It yields none for every polynomial deformation, every
network recursion or the general point-dependent layout interface.

\section{A near-exponential tradeoff for swap-based
fixing}\label{a-near-exponential-tradeoff-for-swap-based-fixing}

\subsection{Construction and theorem}\label{construction-and-theorem}

Fix integers \(r\ge2\) and \(Q\ge1\). Let \[\mathcal H=([r]\times[r])^Q,
\qquad V=\mathbb Q^{\mathcal H},
\qquad D=\dim V=r^{2Q}.\] For \(s\in[Q]\), let \(\Pi_s\) exchange the
two labels at position \(s\). These commuting involutions act on the
coordinate space \(V\) itself. We define \[V_0=\prod_{s=1}^Q\Pi_s,
\qquad W_0=\frac12\sum_{s=1}^Q(I-\Pi_s),
\qquad M_f=V_0f(W_0)\] for \(f\in\mathbb Q[X]\).

Rational linear measurements \(C:V\to\mathbb Q^m\) enter the
construction through one requirement, that a vector with zero
measurements is fixed by \(M_f\):
\[\ker C\subseteq\ker(I-M_f).                 \tag{5.1}\] This is a
property for an application to check. It admits arbitrary rational
measurement rows, in any basis.

\begin{manualstatement}

{{\upshape \textbf{Theorem 5.1 (Measurement, degree, and
enumeration).}}} Let \(r,Q,R,m,n\) be nonnegative integers with
\(r\ge2\), \(Q\ge1\), and \(R\ge1\). Let \(f\in\mathbb Q[X]\) have
degree less than \(R\), let \(\tau\in\mathbb R\) satisfy
\(0\le\tau\le1\), and let \(C:V\to\mathbb Q^m\) be \(\mathbb Q\)-linear
and satisfy (5.1) and
\[m\le\frac{D\tau}{3r}.                       \tag{5.2}\] Then
\[R\ge1+\frac{5\sqrt Q}{24}.                  \tag{5.3}\] If also
\(n\ge2^{16}\) and a real dimension budget \(M\) satisfies
\[12D^2\le M\le3n,                            \tag{5.4}\] then
\(Q\le\tfrac14\log_2 n\), and, with \(q=\lfloor3nR/Q\rfloor\),
\[A(n,q)\ge
2^{\,n\log_2\log_2 n/(4\sqrt{\log_2 n})}.   \tag{5.5}\] For every fixed
real \(\varepsilon>0\), there is a threshold \(N_\varepsilon\),
independent of \(r,Q,R,m,f,\tau,M,C\), such that every choice satisfying
these contracts with \(n\ge N_\varepsilon\) obeys
\[2^{n^{1-\varepsilon}}\le A(n,q)\le3^n.       \tag{5.6}\]

\end{manualstatement}

Here the quantification runs over the operators and measurement maps
themselves. Neither a spectrum nor a degree lower bound enters as a
hypothesis, since the proof establishes both. Its last statement is
uniform over the whole family, and it charges an implementation that
visits every member of the enumeration at least once. A different
computation of the resulting sum is a different question.

\subsection{Exact spectrum and the measurement
requirement}\label{exact-spectrum-and-the-measurement-requirement}

On a single pair space, use the basis consisting of the diagonal vectors
\(e_{aa}\), the symmetric vectors \(e_{ab}+e_{ba}\) for \(a<b\), and the
antisymmetric vectors \(e_{ab}-e_{ba}\) for \(a<b\), so that the swap
has positive and negative eigenspace dimensions
\[s=\frac{r(r+1)}2,\qquad a=\frac{r(r-1)}2.\] Tensoring these bases
produces a basis of \(V\). A tensor with exactly \(j\) antisymmetric
factors has \(W_0\)-eigenvalue \(j\) and \(V_0\)-eigenvalue \((-1)^j\),
with multiplicity \(\binom Qj a^j s^{Q-j}\). Consequently,
\[\operatorname{rank}(I-M_f)
 =\sum_{\substack{0\le j\le Q\\ f(j)\ne(-1)^j}}
   \binom Qj a^j s^{Q-j}
 =D\Pr[f(W)\ne(-1)^W],                     \tag{5.7}\] where
\[W\sim\operatorname{Bin}(Q,p),
\qquad p=\frac{r-1}{2r}\in[1/4,1/2].\]

The commutative-algebra reading of (5.7) appears in \citep[Section
1.2]{T26note}. Here the rational pair and tensor bases, the
multiplicities and the operator identity are all established in the
formal development.

Kernel containment (5.1) and rank--nullity imply
\[\operatorname{rank}(I-M_f)\le\operatorname{rank}C\le m,\] so together
with (5.2) the exact identity (5.7) gives
\[\Pr[f(W)\ne(-1)^W]\le\frac\tau{3r}\le\frac16.  \tag{5.8}\]

In the source construction, selecting every eigendirection outside an
interpolation window is enough to enforce the fixing condition. That
selected count may exceed the true defect rank, since the polynomial may
happen to agree with parity elsewhere. Equation (5.7) counts the defect
itself, so the argument goes through without equating the two
quantities.

\subsection{Why a low measurement budget forces
degree}\label{why-a-low-measurement-budget-forces-degree}

If \(f\) is neither constant \(1\) nor constant \(-1\), its agreement
set with parity on \(\{0,\ldots,Q\}\) has size at most \(2(R-1)\). Each
agreeing point is a root of \(f-1\) or of \(f+1\), and each of those
polynomials has at most \(R-1\) roots.

For \(p\in[1/4,1/2]\) and \(Q\ge1\), the elementary binomial estimate
proved in Appendix A gives \[\max_j\Pr[W=j]\le\frac2{\sqrt Q}.\] Hence
\[\Pr[f(W)=(-1)^W]\le\frac{4(R-1)}{\sqrt Q}.\] Combining this with (5.8)
gives \(5/6\le4(R-1)/\sqrt Q\), which proves (5.3).

Both constant cases need their own argument, and the binomial identity
\(\mathbb E[(-1)^W]=(1-2p)^Q\) supplies it:
\[\Pr[W\text{ odd}]=\frac{1-(1-2p)^Q}{2}\ge\frac14,
\qquad
\Pr[W\text{ even}]\ge\frac12,\] so each constant sign violates (5.8),
and the zero polynomial is already covered by the nonconstant-sign
argument.

Agreement of low-degree polynomials with parity is a classical subject
\citep{ABFR94}. Our \(\sqrt Q\) scale matches interpolation on a central
binomial window, as used in constructions of probabilistic polynomials
\citep{AW15}, and sharp uniform atom bounds are known \citep{BCV16}. We
use an elementary and deliberately loose constant because it gives a
completely checked implication from the measurement contract as stated.
We claim no new general parity lower bound.

\subsection{From degree to the complete
enumeration}\label{from-degree-to-the-complete-enumeration}

Set \(L_2=\log_2 n\). Since \(D=r^{2Q}\ge2^{2Q}\), (5.4) implies
\(2^{4Q}\le n\), so \(Q\le L_2/4\). Combining this with (5.3) and the
floor convention gives \[q\ge\frac{5n}{4\sqrt{L_2}}-1.\] For
\(n\ge2^{16}\), put \(k=\lfloor n/\sqrt{L_2}\rfloor\). Elementary
estimates give \[1\le k\le\min(q,n),
\qquad
\frac{n}{2\sqrt{L_2}}\le k\le\frac n{\sqrt{L_2}}.\] Keeping just the
\(k\)th term of the complete sum and using \(\binom nk\ge(n/k)^k\)
yields \[A(n,q)\ge\binom nk
 \ge\left(\sqrt{L_2}\right)^{n/(2\sqrt{L_2})}
 =2^{\,n\log_2 L_2/(4\sqrt{L_2})},\] which is (5.5). Clipping at \(n\)
matters here, since the argument still applies when \(q>n\). By the
binomial theorem, the complete sum is always at most
\(\sum_{j=0}^n2^j\binom nj=3^n\).

Finally, for every fixed \(\varepsilon>0\) the exponent in (5.5)
eventually exceeds \(n^{1-\varepsilon}\). Indeed
\(4\log_2 n\le n^\varepsilon\) eventually, and once \(L_2\ge16\) we have
\(\sqrt{L_2}\le L_2\) and \(\log_2 L_2\ge1\), which together give the
comparison. Since that threshold depends on \(\varepsilon\) alone, (5.6)
follows. \(\square\)

\subsection{The source measurement
map}\label{the-source-measurement-map}

For the measurements the source itself defines, the fixing condition can
be discharged outright. Let \((b_\alpha)_{\alpha\in\mathcal E}\) be the
orthogonal tensor basis from Section 5.2, and let \(\nu(\alpha)\) count
its antisymmetric factors. For a retained set
\(J\subseteq\{0,\ldots,Q\}\), put
\[\mathcal S_J=\{\alpha:\nu(\alpha)\notin J\},
\qquad m_J=|\mathcal S_J|,\] and define the row map by
\[C_J:V\longrightarrow\mathbb Q^{\mathcal S_J},
\qquad (C_Jv)_\alpha=\sum_{u\in\mathcal H}b_\alpha(u)v(u).\] Ordering
the selected rows identifies its codomain with \(\mathbb Q^{m_J}\).

\begin{manualstatement}

{{\upshape \textbf{Proposition 5.2 (Selected spectral measurements).}}}
Suppose \(f(j)=(-1)^j\) for every \(j\in J\). Then \(C_Jv=0\) implies
\(M_fv=v\). Every matrix entry of \(C_J\) is \(0\), \(1\) or \(-1\),
every row has \(\ell_1\)-norm at most \(D\), and
\[m_J=\sum_{j\notin J}\binom Qj a^j s^{Q-j}
    =D\Pr[W\notin J],\] where the sum ranges over \(0\le j\le Q\), and
\(a,s,W\) are as in Section 5.2. Consequently, the tail condition
\(\Pr[W\notin J]\le\tau/(3r)\) supplies the row budget in Theorem 5.1
for this constructed map, with no fixing-kernel premise left to assume.

\end{manualstatement}

\begin{manualproof}

\textbf{Proof.} Expand \(v=\sum_\alpha\beta_\alpha b_\alpha\).
Orthogonality gives \((C_Jv)_\alpha=\|b_\alpha\|^2\beta_\alpha\). Each
squared norm is positive, so zero measurements force every selected
coefficient to vanish. Every remaining basis vector has
\(\nu(\alpha)\in J\) and \(M_f\)-eigenvalue
\((-1)^{\nu(\alpha)}f(\nu(\alpha))=1\), which proves the fixing
property. Basis entries lie in \(\{0,1,-1\}\), giving the row-height
bound. Counting selected indices by their antisymmetric factors gives
the exact multiplicities and the binomial normalization. \(\square\)

\end{manualproof}

These are precisely the transposed eigenvector rows in the proof of
OpenAI's auxiliary-matrix lemma \citep[Lemma
\texttt{lem:swaps-matrix}]{O26}. In the source, \(J\) is a consecutive
window and \(f\) its degree-less-than-\(|J|\) parity interpolant. Our
formal artifact constructs such an interpolant, proves its values,
degree and uniqueness, and composes it with the selected-row map and
Theorem 5.1; an integer window is handled by retaining its intersection
with the possible eigenlevels. Where the source writes its interpolant
with a Newton binomial formula, the formal application uses the unique
interpolation characterization instead.

This local map \(C_J=C^{\mathcal B}\) is distinct from the formula's
placed measurements, which the source obtains by taking \(h\) block
copies, pulling them back along the variable-repetition map \(\iota\),
and restricting columns to each group:
\[\widetilde C=\operatorname{diag}(C_J,\ldots,C_J),
\qquad \widehat C=\widetilde C\iota,
\qquad C_{t,g}=\widehat C|_{S_{t,g}}.\] Zero rows are retained. In the
bounded-occurrence clause branch, the source has \(hD\le2M\), incidence
and scope bounds \(b=w=3\), and \(12D^2\le M\le3n\). The local tail
budget therefore gives the global row accounting
\[rh\,m_J\le\frac{hD\tau}{3}
          \le\frac{2M\tau}{3}\le2\tau n,\] and the artifact checks this
numerical implication too. Since the placement construction itself comes
from the source, its global row bound on its own would leave the local
tail budget underived. In the source, \(0<\tau\le1/2\); the theorem also
admits the nonnegative boundary.

Its sparse polynomial equation branch \citep[Sections 4.7 and
5.2--5.5]{O26} uses the same local swap operator and measurement rows
over \(\mathbb Q\). On a nontrivial branch with \(n'\) free variables,
the layout instead requires \(2(b+w)D^2\le M\le bn'\) and
\(2Qr^2\le n'\), with \(b\) and \(w\) bounded by powers of \(\log n'\),
and its local budget is \(m_J\le D\tau/(rb)\), where
\(\tau=1/(r\lambda)\) and \(\lambda=\lceil\log_2 n'\rceil\). Thus
\(\tau_* = 3\tau/b\le1\) for sufficiently large \(n'\), and (5.2) holds
with \(\tau_*\) in place of \(\tau\), so the same degree floor (5.3)
follows.

For the enumeration step this branch has its own cutoff
\(q=\lfloor dR/Q\rfloor\), where the source takes
\(d=k(m_{\rm test}+2(\ell-1))\) for degree bound \(k\ge1\),
\(m_{\rm test}\) equations and precision \(\ell=n+1\) on the original
\(n\ge n'\) variables, so \(d\ge2n'\). With the source's choice
\(Q=\Theta(\log n'/\log\log n')\), (5.3) gives \(q=\Omega(n'/\sqrt Q)\),
hence eventually \(q\ge\lfloor n'/\sqrt{\log_2 n'}\rfloor\). Keeping
that term as in Section 5.4 yields
\(\log_2 A(n',q)=\Omega(n'\log\log n'/\sqrt{\log n'})\). So the
near-exponential obstruction applies to the source's complete
enumeration in this branch as well, with its own cutoff and budgets;
this derivation uses the checked degree bound and the source's parameter
choices.

\subsection{What the construction-specific obstruction
establishes}\label{what-the-construction-specific-obstruction-establishes}

What the theorem closes is one escape route. Within the full history
space, with a single univariate fixing polynomial, a measurement map
satisfying (5.1), the row and dimension budgets, and the complete
substitution enumeration, no change of \(f\) or of the rows of \(C\)
brings this enumeration down to \(2^{O(n^{1-\varepsilon})}\) for any
fixed \(\varepsilon>0\). To absorb the constant hidden in that
\(O\)-notation, apply (5.6) with \(\varepsilon/2\).

This chain also shows what pays for the residual interaction. Few
measurements leave a large subspace that the operator must fix. Fixing
that subspace requires agreement with parity on substantial binomial
mass, the agreement requirement forces degree, and degree forces a large
substitution family under the dimension budget. So this is a proved
relation among distinguishable directions, joint algebraic operations,
and the cost of one recovery procedure.

Proposition 5.2 supplies the local measurement application, and Section
5.5 accounts for the distinct budgets of the clause and sparse
polynomial equation branches of \citep{O26}. Everything around them
keeps its source-level analysis: the formula reduction, the schedule and
placement construction, the replacement-form support, the coefficient
bounds and the bit costs. In particular, the present formalization
covers neither the full counting algorithm nor a transfer of the same
contracts to \citep{Ant26}. A new representation, a multivariate fixing
operation, a smaller invariant space or a recovery algorithm that avoids
complete enumeration would each call for its own analysis.

\section{A refined hardness program}\label{a-refined-hardness-program}

Within its specified family, Theorem 5.1 shows the same fixed-power
separation that Weak \#ETH asserts. We read that agreement as evidence
about a mechanism, and nothing more: it establishes no unrestricted
hypothesis. Its value for the broader question lies in the shape of the
proof. Its obstruction is attached to a concrete interaction that a
permitted computation must resolve. That obstruction survives every
change the model already allows. Proposition 5.2 ties the statement to
the local measurements the source uses.

A general structural account would have to justify comparable statements
for a far wider class of computations. Three requirements are
unavoidable. First, task-changing alternatives must be genuine
admissible inputs rather than objects in some convenient ambient space;
Proposition 4.1 shows why. Second, any transfer across representations
must account for the information supplied and for its acquisition,
instead of treating an aggregate operation as a point query; Proposition
3.4 gives one sufficient condition, and no classification. Third,
computation must be charged at the level of evaluation and recovery as
performed. That level includes construction, arithmetic precision,
reuse, and the selection of intermediate objects. The network algorithms
show why certificate size alone cannot supply these costs.

This leaves a concrete way to investigate the successor hypothesis. For
a proposed fixed-power improvement, locate the joint operation that
produces the gain. Identify the interaction it leaves unresolved. Then
prove a tradeoff that includes the operation's full contract. Section
4.3 illustrates three realizations of this: graph structure and recovery
exactness, rewrite-and-merge geometry, and support-constrained rank. Our
finite spectral analysis adds a checked instance with a source-defined
measurement map. An obstruction may concern rank, degree, support,
representation size or the cost of a remaining sum. A new algorithm may
instead show that the proposed obstruction was tied to too narrow a
family. Either way the structural claim becomes more precise.

This program is more modest than the slogan that harder problems require
more work. It is no monotonicity theorem in one syntactic size
parameter. Easy instances may have long descriptions. Short descriptions
may encode expensive tasks. The relevant distinctions are relative to
the requested answer and to the operations available for computing it.
Nor does the program hold that every change of representation preserves
hardness. The reported ETH algorithms matter precisely because a
different organization can lower computational cost.

Our quantitative destination remains \#SH rather than the original ETH
boundary. Public reduction theory \citep{SH26} already explains how
fixed-power lower bounds propagate when reduction costs and size blowups
are controlled. The complementary task pursued here is to explain and
prove the internal obstruction for one counting mechanism. Extending
such arguments to arbitrary exact-counting algorithms is open. What the
present results supply is an independently stated, formally verified
case in which the refined near-exponential scale persists.

\section{Formalization and
reproducibility}\label{formalization-and-reproducibility}

Our Lean 4 development uses Mathlib \citep{Lean4, Mathlib20}. Its public
entry point is \texttt{ETHResponse.lean}. Its selected import closure
contains the task-certificate and transcript results, the width-two CNF
construction, operational bit reading, finite observation simulation,
rational Hermite recovery, the swap basis and spectrum, the
measurement-rank implication, and both the finite and the eventual
enumeration bounds.

In the formalization, the swap operator acts on rational functions on
the history product itself. The spectral identity with its binomial
multiplicities is proved from a rational basis rather than assumed
numerically. Its general measurement endpoint quantifies over a rational
linear map satisfying the fixing-kernel condition. An application module
constructs the transpose rows the source selects. It proves their
relation to dual coordinates, their entry and height bounds, their row
count and their kernel property. It composes them with the finite and
eventual bounds. Separate interpolation and window results discharge the
parity-agreement premise for the constructed interpolant. Global
placement matrices lie outside this formal slice.

\texttt{ETHResponse/\allowbreak{}Audit.\allowbreak{}lean} checks the
transitive axiom dependencies of every loaded declaration in the
project's mathematical namespaces, private helpers included. It permits
only propositional extensionality, classical choice and quotient
soundness. An export records source hashes, pinned dependency revisions,
build results and the freshly measured declaration count.
\texttt{CORRESPONDENCE.md} maps the paper's mathematical statements to
named theorem endpoints and their application hypotheses.

So the formalization establishes the mathematical implications stated
here, and those alone. \#SH itself, its validity, an end-to-end
verification of either external counting algorithm, and a general time
lower bound from the point-access example all lie outside it. Each
accompanying TeX source archive ships the selected Lean project in
\texttt{anc/\allowbreak{}lean\_\allowbreak{}proofs/\allowbreak{}}, with
its correspondence document, toolchain and dependency lockfile.
\texttt{SOURCE\_MANIFEST.json} identifies the proof snapshot by file
hashes. The build and axiom receipts bind their results to that
snapshot.

\appendix

\section{An elementary binomial atom
bound}\label{an-elementary-binomial-atom-bound}

We include the estimate used in Section 5 to make the numerical degree
constant independent of an imported asymptotic approximation. Let
\(b_{m,j}=2^{-m}\binom mj\) be the fair-coin mass. The central-binomial
recurrence gives \[b_{m,j}^2\le\frac1{m+1}
\quad\text{for every }m,j\ge0.\] Put \(c_k=4^{-k}\binom{2k}k\). The
recurrence \((2k+2)c_{k+1}=(2k+1)c_k\) gives \((2k+1)c_k^2\le1\) by
induction. The maximum mass is \(c_k\) in row \(2k\) and \(c_{k+1}\) in
row \(2k+1\), proving the bound.

For \(0<p\le1/2\), let \(N\sim\operatorname{Bin}(Q,2p)\). Thinning by
independent fair coins gives
\[\Pr[\operatorname{Bin}(Q,p)=j]=\mathbb E[b_{N,j}].\] Cauchy--Schwarz
and the fair-coin bound imply \[\Pr[\operatorname{Bin}(Q,p)=j]^2
 \le\mathbb E\!\left[\frac1{N+1}\right]
 \le\frac1{2p(Q+1)}.\] For \(p\ge1/4\) and \(Q\ge1\), this is at most
\(2/(Q+1)\) and hence yields the convenient bound \(2/\sqrt Q\) used
above. The Lean proof establishes the finite mixture identity
algebraically and proves the inverse-moment upper bound; no
probabilistic limit theorem is required.

\section{Principal formal endpoints}\label{principal-formal-endpoints}

The table gives short names; each theorem is in the
\texttt{ExactSolution} namespace. The source correspondence file
includes the supporting modules and exact application contracts.

{\def\LTcaptype{none} % do not increment counter
\begin{longtable}[]{@{}
  >{\raggedright\arraybackslash}p{(\linewidth - 2\tabcolsep) * \real{0.4822}}
  >{\raggedright\arraybackslash}p{(\linewidth - 2\tabcolsep) * \real{0.5178}}@{}}
\toprule\noalign{}
\begin{minipage}[b]{\linewidth}\raggedright
Mathematical statement
\end{minipage} & \begin{minipage}[b]{\linewidth}\raggedright
Lean theorem
\end{minipage} \\
\midrule\noalign{}
\endhead
\bottomrule\noalign{}
\endlastfoot
Completed transcript is a certificate &
\texttt{transcript\_certificate} \\
Transcript query lower bound &
\texttt{transcript\_\allowbreak{}query\_\allowbreak{}lower\_\allowbreak{}bound} \\
Width-two CNF count separation &
\texttt{SingletonCNF.\allowbreak{}singleton\_\allowbreak{}count\_\allowbreak{}complexity} \\
Explicit-input acquisition ceiling & \texttt{read\_all\_operational} \\
Finite weighted observation simulation &
\texttt{weighted\_\allowbreak{}observation\_\allowbreak{}simulation} \\
Ambient Hermite threshold &
\texttt{hermite\_\allowbreak{}value\_\allowbreak{}at\_\allowbreak{}zero\_\allowbreak{}determined\_\allowbreak{}iff} \\
Actual swap defect rank &
\texttt{SwapSpectrum.\allowbreak{}swap\_\allowbreak{}rank\_\allowbreak{}eq\_\allowbreak{}failure\_\allowbreak{}mass} \\
Measurement-to-degree implication &
\texttt{MeasurementRank.\allowbreak{}swap\_\allowbreak{}measurement\_\allowbreak{}degree\_\allowbreak{}floor} \\
Finite clipped-enumeration floor &
\texttt{MeasurementRank.\allowbreak{}swap\_\allowbreak{}measurement\_\allowbreak{}enumeration\_\allowbreak{}floor} \\
Uniform eventual near-exponential floor &
\texttt{swap\_\allowbreak{}measurement\_\allowbreak{}near\_\allowbreak{}exponential} \\
Complete-enumeration upper bound &
\texttt{enumerationCount\_\allowbreak{}le\_\allowbreak{}three\_\allowbreak{}pow} \\
Constructed physical measurement kernel &
\texttt{SelectedMeasuremen\allowbreak{}ts.\allowbreak{}measurement\_\allowbreak{}fixing\_\allowbreak{}kernel} \\
Integer measurement-row height &
\texttt{SelectedMeasuremen\allowbreak{}ts.\allowbreak{}integerMeasurement\allowbreak{}Matrix\_\allowbreak{}l1\_\allowbreak{}le} \\
Exact selected-row tail count &
\texttt{SelectedMeasuremen\allowbreak{}ts.\allowbreak{}rowCount\_\allowbreak{}eq\_\allowbreak{}dimension\_\allowbreak{}mul\_\allowbreak{}tail} \\
Interpolant uniqueness &
\texttt{SelectedMeasuremen\allowbreak{}ts.\allowbreak{}eq\_\allowbreak{}parityInterpolant} \\
Constructed interval enumeration &
\texttt{SelectedMeasuremen\allowbreak{}ts.\allowbreak{}interval\_\allowbreak{}measurement\_\allowbreak{}enumeration\_\allowbreak{}floor} \\
Uniform interval application &
\texttt{SelectedMeasuremen\allowbreak{}ts.\allowbreak{}interval\_\allowbreak{}measurement\_\allowbreak{}near\_\allowbreak{}exponential} \\
Placed-row numerical accounting &
\texttt{SelectedMeasuremen\allowbreak{}ts.\allowbreak{}placed\_\allowbreak{}row\_\allowbreak{}count\_\allowbreak{}bound} \\
\end{longtable}
}
\bibliographystyle{plainnat}
\bibliography{references}
\end{document}
